\begin{enumerate}[label=\textbf{\ansnum.\arabic*}]
  \item Example: (1) batteries charged and numbered, (2) spare parts, (3) the emergency stop works on every robot, (4) jerseys and numbers, (5) the software version is frozen and tagged, (6) the GameController test passes, (7) IP addresses match the assigned range, (8) camera calibration in the light of the hall, (9) names of handlers, captain and referee, (10) the rulebook is printed. In the lab: 1, 2, 3, 4, 5, 6. Only on the field: 7 (IP range and network load), 8 (the real light), partly 6 (the real access point).
  \item Less than two hours before the game: this counts as a forfeit. The opponent plays against an empty field and the own goals of the forfeiting team are not counted; a qualification penalty is recorded. (With at least two hours of notice, in a round-robin game the same happens; in a knock-out game the team is replaced by the next best qualified team.)
  \item The penalty timer does \emph{not} run during maintenance: the repair takes $30$\,s of maintenance, so the robot has still $45$\,s to serve after the repair, and returns $75$\,s after the penalty started (if the repair is done at the beginning).
  \item In the draft rules a robot that is inactive for $10$\,s is incapable, and a fallen robot must try to get up within $20$\,s. If the robot lies still it is removed after about $10$\,s of inactivity (at the latest after $20$\,s). The team then loses it for $45$\,s plus the time before the removal: between $10+45=55$\,s and $20+45=65$\,s of robot time, plus the walk back into the game. This is why a fast and reliable get-up is part of the game rules.
  \item Needed: state, set play and kicking team, the time left, the score and its own penalty status (secs till unpenalised), and the message budget. The robot must \emph{re-estimate} its position and orientation (it was moved by hand, so its old pose is wrong) and the ball; the particle filter of Chapter~\ref{ch:ctl5} must be reset after a pick-up (``kidnapped robot'').
  \item Record the time at which the operator presses the button (the GameController log has it) and the time at which the robot reports its safe pose (a log in the robot, or a camera at $120$ frames per second). The difference is the reaction time; repeat $20$ times and report the worst case, not the mean. Include the network delay and the control period of the robot.
  \item The rule allows one packet per second per robot to \emph{one} device of the team on the \emph{wired} network; a $30$\,kB image every $5$\,s is only $0.2$ packets/s, so the count is no problem, but the size of a single UDP packet is limited to about $64$\,kB, and it is the wired network (you cannot use the wireless access point). To get more data out of a game, \emph{log on the robot} and read the log after the game.
  \item Example: Hardware: spare batteries, spare servo or leg parts, a charger, a switch and cables for the wired network, the robot's remote emergency button. Software: a laptop with the tagged software and a copy of the GameController, the calibration tools, a USB stick with the complete code (offline), the log analysis tool, the rulebook.
  \item Requirement: ``After a fall (the IMU detects a lying robot) the get-up behaviour starts within $1$\,s, and the robot is standing within $20$\,s in $95\%$ of the falls, with at most three attempts.'' Test: push the robot over $100$ times (forward, backward, side), measure the time from the fall to standing; pass criterion: $\ge95$ successes, maximum $20$\,s.
  \item A short program: read the messages, accumulate the time per state with $0.5$\,s per message, count the changes of \emph{stopped} from $0$ to $1$, and take the minimum of the budget. For the test log: $20\,\text{min}=2400$ messages, the budget goes from $12000$ to $12000-1.6\cdot5\cdot1200=2400$ (if the budget is only used during playing), so the report must show $2400$ as the lowest value, $1200$\,s in the state \emph{playing} (as generated), and the number of stops that you put in. (A sketch; no program of this problem was run.)
\end{enumerate}
